\chapter{Mise en œuvre II : jouer sans savoir évaluer}

Deuxième application complète, dans un tout autre monde : l'intelligence
artificielle d'un jeu de plateau. Le principe est identique, les gains sont
spectaculaires, et la mise en œuvre tient en une page.

\section{Le problème des méthodes classiques}

Un moteur de jeu traditionnel (Minimax, Negamax, alpha-bêta) explore l'arbre
des coups possibles et, à une certaine profondeur, appelle une \textbf{fonction
d'évaluation} : « cette position vaut $+3$ pour les blancs ».

Cette fonction est le talon d'Achille. Aux échecs, on sait la fabriquer
(matériel, structure de pions, sécurité du roi — des siècles de théorie). Mais
sur un jeu nouveau, sur le Go, ou sur votre propre jeu de plateau,
\textbf{personne ne sait dire si une position à mi-partie est bonne}. Et une
mauvaise fonction d'évaluation produit une IA confiante et stupide.

\begin{nkretenir}[L'idée qui contourne le problème]
Ne jugeons pas la position. \textbf{Jouons la partie jusqu'au bout}, au hasard,
et regardons qui gagne. À la fin d'une partie, il n'y a plus rien à évaluer :
le résultat est un fait.

« Je ne sais pas si ce coup est bon. Alors depuis ce coup, je joue mille
parties complètement au hasard jusqu'à la fin. Si j'en gagne sept cents, ce
coup est probablement bon. »
\end{nkretenir}

C'est du Monte-Carlo pur : on estime une valeur inconnue (la qualité d'un coup)
par la moyenne de tirages aléatoires (des parties jouées au hasard).

\section{Le raffinement : ne pas répartir équitablement}

Jouer mille parties par coup gaspille : la plupart des coups sont manifestement
mauvais après cinquante parties, et continuer à les explorer prive les coups
prometteurs de temps de calcul.

MCTS (\emph{Monte Carlo Tree Search}) résout ce dilemme avec une formule qui
arbitre entre \textbf{exploiter} ce qui marche et \textbf{explorer} ce qu'on
connaît mal :

\[ UCB1(i) \;=\; \underbrace{\frac{w_i}{n_i}}_{\text{ce que ça rapporte}}
   \;+\; \underbrace{c\,\sqrt{\frac{\ln N}{n_i}}}_{\text{curiosité}} \]

où $w_i$ est le total des scores obtenus par ce coup, $n_i$ son nombre de
visites, et $N$ le nombre de visites du nœud parent.

\begin{itemize}
  \item le premier terme est le \textbf{taux de victoire} : il pousse vers ce
        qui a marché ;
  \item le second \textbf{grossit quand $n_i$ est petit} : il force à essayer
        les coups délaissés. Le $\ln N$ au numérateur fait que cette curiosité
        décroît lentement à mesure qu'on connaît mieux le nœud ;
  \item $c$ règle l'équilibre. \textbf{$c = \sqrt{2}$} si les scores sont dans
        $[0,1]$ : c'est la valeur qui découle de la théorie des bandits
        manchots. Plus grand rend l'IA curieuse, plus petit la rend gourmande.
\end{itemize}

\begin{nkretenir}[MCTS est de l'échantillonnage d'importance]
C'est exactement l'idée du chapitre 4, appliquée à un arbre : \textbf{tirer
davantage là où ça compte}. UCB1 est la densité $p$, et l'arbre se construit
lui-même autour des branches qui le méritent.
\end{nkretenir}

\section{Les quatre phases, répétées en boucle}

\begin{enumerate}
  \item \textbf{Sélection} — depuis la racine, descendre en choisissant à
        chaque nœud l'enfant de plus fort UCB1, jusqu'à un nœud dont tous les
        coups n'ont pas encore été essayés ;
  \item \textbf{Expansion} — créer un nouvel enfant pour l'un de ces coups
        jamais tentés ;
  \item \textbf{Simulation} — depuis ce nouveau nœud, jouer \textbf{au hasard}
        jusqu'à la fin de la partie. Résultat : victoire, défaite ou nul ;
  \item \textbf{Rétropropagation} — remonter jusqu'à la racine en ajoutant une
        visite et le score à chaque nœud traversé.
\end{enumerate}

\begin{nkpiege}[L'erreur nº1, celle qui rend l'IA suicidaire]
À la remontée, le score doit être \textbf{inversé à chaque changement de
joueur}. Une victoire pour vous est une défaite pour l'adversaire.

Si vous remontez le même chiffre partout, votre IA choisit les coups qui font
gagner \emph{l'autre} — et elle le fait avec une parfaite assurance
statistique. Le symptôme est déroutant : plus vous lui donnez de temps de
calcul, plus elle joue mal.
\end{nkpiege}

\section{Le code complet}

\begin{nkcode}{MCTS -- structure de noeud et boucle principale}
struct NkMctsNode {
        NkMctsNode* parent = nullptr;
        NkVector<NkMctsNode*> enfants;
        NkVector<NkCoup> coupsNonEssayes;  // vide = noeud pleinement etendu
        NkCoup coupJoue;                    // le coup qui mene ICI
        nk_uint32 visites = 0;
        float32 score = 0.f;                // du point de vue de `joueur`
        nk_int32 joueur = 0;                // qui a joue `coupJoue`
};

NkCoup MeilleurCoup(const NkPlateau& depart, nk_uint32 budget) {
    NkMctsNode* racine = CreerNoeud(nullptr, depart);

    for (nk_uint32 iter = 0; iter < budget; ++iter) {
        NkMctsNode* noeud = racine;
        NkPlateau etat = depart;             // copie de travail

        // (1) SELECTION : descendre tant que le noeud est etendu.
        while (noeud->coupsNonEssayes.Empty() && !noeud->enfants.Empty()) {
            noeud = MeilleurEnfantUCB(noeud, 1.41421356f);   // c = sqrt(2)
            etat.Jouer(noeud->coupJoue);
        }

        // (2) EXPANSION : essayer un coup encore jamais joue.
        if (!noeud->coupsNonEssayes.Empty()) {
            const nk_uint32 k =
                math::NkRand.NextUInt32(noeud->coupsNonEssayes.Size());
            const NkCoup coup = noeud->coupsNonEssayes[k];
            noeud->coupsNonEssayes.Erase(k);
            etat.Jouer(coup);
            noeud = AjouterEnfant(noeud, coup, etat.JoueurCourant());
        }

        // (3) SIMULATION : jouer au hasard jusqu'au bout.
        NkPlateau sim = etat;
        NkVector<NkCoup> coups;
        while (!sim.EstTerminee()) {
            coups.Clear();
            sim.CoupsLegaux(coups);
            if (coups.Empty())
                break;
            sim.Jouer(coups[math::NkRand.NextUInt32(coups.Size())]);
        }
        const nk_int32 gagnant = sim.Gagnant();

        // (4) RETROPROPAGATION -- le score DEPEND du joueur de chaque noeud.
        while (noeud != nullptr) {
            noeud->visites += 1;
            if (gagnant == 0)
                noeud->score += 0.5f;                        // partie nulle
            else
                noeud->score += (gagnant == noeud->joueur) ? 1.f : 0.f;
            noeud = noeud->parent;
        }
    }

    // Le coup rendu est le PLUS VISITE, pas le mieux note : le nombre de
    // visites est une statistique bien plus stable qu'un taux de victoire
    // calcule sur trois parties. Un coup vu deux fois avec deux victoires
    // affiche 100 % -- et ne veut rien dire.
    return EnfantLePlusVisite(racine)->coupJoue;
}
\end{nkcode}

\begin{nkcode}{UCB1 -- le choix de l'enfant}
NkMctsNode* MeilleurEnfantUCB(NkMctsNode* n, float32 c) {
    NkMctsNode* meilleur = nullptr;
    float32 meilleureValeur = -1e30f;
    const float32 logN = NkMath::Log((float32)n->visites);

    for (nk_size i = 0; i < n->enfants.Size(); ++i) {
        NkMctsNode* e = n->enfants[i];

        // Un enfant jamais visite est prioritaire ABSOLU : sa formule
        // diviserait par zero, et la theorie veut qu'on l'essaie d'abord.
        if (e->visites == 0)
            return e;

        const float32 exploitation = e->score / (float32)e->visites;
        const float32 exploration  = c * NkMath::Sqrt(logN / (float32)e->visites);
        const float32 v = exploitation + exploration;
        if (v > meilleureValeur) {
            meilleureValeur = v;
            meilleur = e;
        }
    }
    return meilleur;
}
\end{nkcode}

\section{Ce que MCTS vous donne gratuitement}

\begin{itemize}
  \item \textbf{aucune connaissance du jeu n'est requise} : il suffit de savoir
        énumérer les coups légaux et reconnaître une fin de partie. Le même
        code marche sur un jeu que vous venez d'inventer ;
  \item \textbf{interruptible} : arrêtez-le quand vous voulez, il rend son
        meilleur coup actuel. Une IA qui doit répondre en 200 ms et une qui a
        10 secondes, c'est le même code avec un budget différent ;
  \item \textbf{il s'améliore tout seul} avec le temps de calcul, sans réglage ;
  \item \textbf{il se parallélise} : plusieurs fils peuvent explorer l'arbre en
        même temps (avec les précautions du chapitre 7).
\end{itemize}

\begin{nkpiege}[Ce que MCTS fait mal]
Les positions à \textbf{coup unique salvateur} — celles où un seul coup sauve
la partie et tous les autres perdent. Le hasard des simulations le trouve
rarement, et la moyenne le noie. C'est pourquoi les meilleurs programmes
combinent MCTS avec une recherche exhaustive à faible profondeur, ou avec un
réseau de neurones qui oriente les simulations (c'est la recette d'AlphaGo :
MCTS + un réseau pour remplacer le tirage uniforme).
\end{nkpiege}

\begin{nkexo}[Le sentir sur un morpion]
Implémentez MCTS pour le tic-tac-toe, jeu assez petit pour qu'on connaisse la
vérité : bien joué, il est toujours nul.

Faites varier le budget : 100, 1000, 10\,000 itérations. À 10\,000, votre IA ne
doit \textbf{jamais} perdre. Puis inversez volontairement le signe de la
rétropropagation : elle deviendra suicidaire — et vous reconnaîtrez ce symptôme
à vie.
\end{nkexo}
